\bookchapter{الفصل 40: هندسة المتطلبات — تحويل الأفكار إلى تعليمات يستطيع الذكاء الاصطناعي تنفيذها}

\chapterintro{إذا كانت هندسة السياق تحل مشكلة «جعل الذكاء الاصطناعي يفهم الخلفية»، فإن هندسة المتطلبات تحل مشكلة «جعل الذكاء الاصطناعي يعرف ما المطلوب فعله تحديدًا هذه المرة». وهذه هي المرحلة التي تتجلى فيها الفوارق بين الناس أكثر من غيرها في البرمجة بالذكاء الاصطناعي — فبالأدوات نفسها والنماذج نفسها، ينجح بعضهم في مشاريع مبهرة بينما لا ينتج آخرون سوى أشياء بسيطة عادية. والغالب أن الفرق لا يكمن في القدرة التقنية، بل في القدرة على التعبير عن المتطلبات.}

\booksection{أولًا: من الفكرة الغامضة إلى المتطلبات المنظمة}

يبدأ كل مشروع بفكرة غامضة — «أريد عمل تطبيق لتسجيل المصروفات»، «أريد عمل موقع شخصي»، «أريد عمل أداة لقائمة المهام». هذه الأفكار لا مشكلة فيها في ذاتها، لكن إن ألقيتها مباشرة إلى الذكاء الاصطناعي فإن جودة المخرجات ستعتمد كليًا على الحظ. هدف هندسة المتطلبات هو تحويل هذه الفكرة الغامضة إلى متطلبات منظمة يستطيع الذكاء الاصطناعي تنفيذها بدقة.

والخبر الجيد: المتطلبات المنظمة لا تتطلب منك كتابة وثيقة متطلبات (PRD) من عشرات الصفحات، بل يكفي أن تجيب عن خمسة أسئلة نسميها «العناصر الخمسة للمتطلبات»:

\begin{center}
\begin{tabularx}{\linewidth}{|X|X|X|}
\hline
\textbf{العنصر} & \textbf{السؤال} & \textbf{مثال: تطبيق تسجيل المصروفات} \tabularnewline
\hline
Who & لمن سيستخدمه؟ & استخدام شخصي لتسجيل الإيرادات والمصروفات اليومية \tabularnewline
What & ماذا نبني؟ & تطبيق يسجل كل عملية إيراد ومصروف، ويجمعها شهريًا، ويعرضها بصريًا \tabularnewline
Why & لماذا نبنيه؟ & لأفهم هيكل إنفاقي وأكتشف النفقات غير الضرورية \tabularnewline
How & كيف نبنيه تقريبًا؟ & الواجهة الأمامية بـ React، والبيانات تُخزَّن محليًا في المتصفح، مع تأجيل الخادم الخلفي وتسجيل الدخول \tabularnewline
Constraint & ما القيود؟ & تطوير فردي، والإنجاز خلال أسبوعين، دون حاجة إلى النشر والإطلاق \tabularnewline
\hline
\end{tabularx}
\end{center}

\begin{center}\small\color{gray}العناصر الخمسة للمتطلبات مع الأمثلة\end{center}

لا حاجة إلى كتابة هذه العناصر الخمسة في وثيقة رسمية، بل يكفي حتى استعراضها سريعًا في ذهنك. والمهم هو اكتساب عادة «التفكير بوضوح قبل الكلام». وستجد أنه حين تستطيع الإجابة عن هذه الأسئلة الخمسة بوضوح، يصبح وصفك للمتطلبات منظمًا تلقائيًا.

\booksection{ثانيًا: فن تقسيم المتطلبات إلى طبقات}

في المشاريع الأكبر لا يكفي وصف متطلبات واحد، بل تحتاج إلى تقسيم المتطلبات الكبيرة إلى متطلبات فرعية صغيرة قابلة للتحقق منها كل على حدة. تسمى هذه العملية «تقسيم المتطلبات»، ومبدأها الأساسي أن كل متطلب فرعي يجب أن يمكن إنجازه والتحقق منه وتسليمه بشكل مستقل.

بأخذ مثال «عمل تطبيق شخصي لتسجيل المصروفات»، فإحدى الطرق المعقولة لتقسيمه هي:

\begin{enumerate}
  \item تنفيذ وظيفة أساسية لتسجيل الإيرادات والمصروفات (إمكانية إضافة السجلات وعرضها وحذفها)؛
  \item إضافة وظيفة التصنيفات (وسم كل سجل بفئة مثل المطاعم أو المواصلات)؛
  \item تنفيذ وظيفة الجمع الشهري (إحصاء الإيرادات والمصروفات شهرًا بشهر)؛
  \item إضافة العرض البصري للبيانات (إظهار هيكل الإنفاق برسوم بيانية).
\end{enumerate}

كل خطوة هي إضافة وظيفية يمكن تشغيلها والتحقق منها بشكل مستقل.

القاعدة الذهبية لتقسيم المتطلبات: بعد إتمام كل خطوة يجب أن يكون المشروع "شيئًا قابلًا للاستخدام"، وإن كانت وظائفه بدائية. لا تترك المشروع أبدًا في حالة "كل شيء نصف منجز ولا شيء يعمل".

لهذا التقسيم ميزتان. أولهما أنه يُبقيك ترى التقدم دائمًا — فمع إتمام كل خطوة يكتسب تطبيقك وظيفة جديدة، وهذه التغذية الراجعة الإيجابية بالغة الأهمية في المشاريع طويلة الأمد (فالجميع يحتاج إلى شعور بالإنجاز). وثانيهما أنه يجعل كل مهمة موكولة إلى الذكاء الاصطناعي صغيرة بما يكفي، ما يقلل خطر «سوء فهم الذكاء الاصطناعي» — فكلما صغرت المهمة كان فهمه أدق وارتفعت جودة المخرجات.

\booksection{ثالثًا: قالب كتابة وثيقة المتطلبات (PRD) والتطبيق العملي}

في المشاريع الأكبر (مثل تطبيق متكامل فيه واجهة أمامية وخلفية وعدة وحدات وظيفية)، قد لا تكفي العناصر الخمسة وحدها. عندها تحتاج إلى كتابة وثيقة متطلبات مشروع (PRD) أكثر رسمية. ونقطة يسهل سوء فهمها: وثيقة PRD ليست مكتوبة للذكاء الاصطناعي — بل هي مكتوبة لك أنت لتساعدك على ترتيب أفكارك، وبعد ذلك فقط تستطيع إعطاء الذكاء الاصطناعي تعليمات أكثر دقة.

وثيقة PRD الجيدة لا تحتاج إلى أن تكون طويلة، لكن ينبغي أن تتضمن ستة أقسام: نظرة عامة على المشروع (فقرة توضح ماذا نبني)، وقائمة الوظائف (ما هي كل وظيفة)، والحل التقني (أي حزمة تقنية وأي بنية تقريبية)، ونموذج البيانات (هياكل البيانات الأساسية)، ووصف الواجهات (شكل الصفحات الرئيسية)، وخطة التطوير (على كم خطوة سيُنجز المشروع).

\begin{lstlisting}
# وثيقة متطلبات تطبيق تتبع المصروفات الشخصية
## وصف المشروع
تطبيق خفيف لتتبع المصروفات الشخصية يدعم تسجيل الإيرادات والمصروفات والتصنيف والملخص الشهري وعرض البيانات بيانيًا.
## قائمة الوظائف
1. تسجيل الحركات: المبلغ، النوع (إيراد/مصروف)، التصنيف، الملاحظة، التاريخ
2. إدارة التصنيفات: طعام، مواصلات، تسوق، ترفيه، أخرى (معدة مسبقًا)
3. الملخص الشهري: إجمالي الإيرادات وإجمالي المصروفات والرصيد لكل شهر
4. عرض المصروفات: مخطط دائري يوضح نسبة كل تصنيف
5. تصدير البيانات: تصدير إلى CSV
## الحل التقني
- الواجهة الأمامية: React + TypeScript + Tailwind CSS
- التخزين: localStorage (بدون خادم)
- الرسوم: Chart.js
- البناء: Vite
## نموذج البيانات
Record {
  id: string
  amount: number
  type: 'income' | 'expense'
  category: string
  note: string
  date: string  // YYYY-MM-DD
}
## خطة التطوير
Step 1: التسجيل الأساسي للحركات (CRUD)
Step 2: إدارة التصنيفات
Step 3: الملخص الشهري
Step 4: عرض البيانات
Step 5: تصدير البيانات
\end{lstlisting}

لاحظ أسلوب كتابة هذه الوثيقة: قوائم موجزة وتنسيق منظم بدلًا من الإطالة والحشو. كل وظيفة لها تعريف واضح، ويُوصف نموذج البيانات بأنواع بسيطة. هذا الأسلوب يجعل الوثيقة واضحة لك أنت، وعندما تقدمها إلى الذكاء الاصطناعي يستطيع هو الآخر فهم كل نقطة متطلب بدقة.

\booksection{رابعًا: تطبيق عملي: محاكاة كاملة من الفكرة إلى الخطة}

الكلام النظري وحده لا متعة فيه، فلنمشي معًا عبر عملية هندسة المتطلبات كاملة بمثال متكامل. لنفترض أن فكرتك هي: «أريد عمل متتبع للتعلم يرصد ما أتعلمه كل يوم».

الخطوة الأولى: استعرض العناصر الخمسة للمتطلبات.

\begin{itemize}
  \item Who: استخدام شخصي.
  \item What: أداة تسجل محتوى التعلم كل يوم وتتتبع مدة التعلم وتراجع تقدم التعلم.
  \item Why: لأنني أريد اكتساب عادة التعلم المستمر ورؤية ما أراكمه بنفسي.
  \item How: أصنع أولًا نسخة ويب، وتُخزَّن البيانات محليًا في المتصفح.
  \item Constraint: تطوير بلا خبرة سابقة، وإنجاز الوظائف الأساسية خلال أسبوع.
\end{itemize}

الخطوة الثانية: قسِّم المتطلبات.

\begin{enumerate}
  \item Step 1: تنفيذ وظيفة «تسجيل عناصر التعلم» (التاريخ والمحتوى والمدة)؛
  \item Step 2: إضافة «إحصاء مدة التعلم» (إجمالي اليوم وإجمالي الأسبوع)؛
  \item Step 3: تنفيذ عرض «تقويم التعلم» (يُعرض مثل مخطط المساهمات في GitHub)؛
  \item Step 4: إضافة وظيفة «مراجعة التعلم» (عرض عشوائي لسجل التعلم في يوم ماضٍ للمساعدة على المراجعة).
\end{enumerate}

الخطوة الثالثة: اكتب وثيقة PRD موجزة. لا داعي للإطالة — فالتقسيم أعلاه بذاته يشكل جوهر الوثيقة. رتبها في ملف .md وضعه في المجلد الجذري للمشروع ليكون مرجعًا مشتركًا لك وللذكاء الاصطناعي.

الخطوة الرابعة: ابدأ تطوير الخطوة الأولى. ويمكن أن تكون التعليمة الموجهة إلى الذكاء الاصطناعي كما يلي:

\begin{lstlisting}
أحتاج الوظيفة الأولى لمتتبع تعلم: تسجيل جلسات الدراسة.
يمكن للمستخدم إضافة سجل يتضمن التاريخ ومحتوى الدراسة والمدة (بالدقائق).
تُخزَّن البيانات في localStorage.
الواجهة الأمامية بـ React + TypeScript.
أعطني التنفيذ الكامل.
\end{lstlisting}

لاحظ سمات هذه التعليمة: فهي تصف هدف هذه الخطوة فقط (وظيفة التسجيل) دون خلطها بوظائف أخرى، وتحدد القيود التقنية (React وTypeScript وlocalStorage)، وتقدم معلومات كافية ليعرف الذكاء الاصطناعي «ماذا يفعل» و«كيف يفعل»، مع ترك التنفيذ التفصيلي له. هذه هي تعليمة برمجة عالية الجودة مرت بمرحلة هندسة المتطلبات.

\begin{tcbpractice}{هيّا نطبق: اكتب بطاقة العناصر الخمسة لمشروعك المفضل}
اختر مشروعًا صغيرًا طالما يدور في خاطرك، وأمضِ 5 دقائق في تعبئة بطاقة واحدة: Who / What / Why / How / Constraint، بجملة واحدة لكل عنصر. ثم قسّمها إلى 3-5 خطوات «كل خطوة منها تُنتج شيئًا قابلًا للاستخدام». وستكتشف: بعد تدوين الفكرة تتحول من «حلم نهاري» إلى «قائمة مهام».

\end{tcbpractice}

\begin{tcbsummary}{خلاصة الفصل}
\begin{itemize}
  \item العناصر الخمسة للمتطلبات: Who / What / Why / How / Constraint — فكّر بوضوح قبل أن تتكلم
  \item القاعدة الذهبية لتقسيم المتطلبات: بعد إتمام كل خطوة يبقى المشروع «شيئًا قابلًا للاستخدام»
  \item وثيقة PRD تُكتب لنفسك: نظرة عامة، قائمة الوظائف، الحل التقني، نموذج البيانات، الواجهات، الخطة
  \item التعليمة الجيدة = وصف الخطوة الحالية فقط + قيود واضحة + قول «ماذا نفعل» لا «كيف نفعل»
  \item قالب المحاكاة الكامل: العناصر الخمسة التقسيم وثيقة PRD إعطاء التعليمة، في أربع خطوات
\end{itemize}

\end{tcbsummary}
